\documentclass[main.tex]{subfiles}
\begin{document}

Có một câu hỏi mà tôi đã tự hỏi mình rất nhiều lần trong suốt hành trình viết quyển sách này: ``Nếu một ngày nào đó, bạn tỉnh dậy mà không nhớ mình là ai, liệu bạn có đủ can đảm để tìm lại chính mình không?''

Đó chính là khởi điểm của mọi thứ.

Hikawa Kuon - hay đúng hơn, ``Haragen Kyuen'' như cậu ấy từng được gọi - đã tỉnh dậy trên một hành lang trường học vắng lặng, mất sạch ký ức, bên cạnh là một cô gái tự xưng là em gái song sinh mà cậu chẳng nhớ nổi. Cái tên giả, bộ đồng phục lạ, và một ngôi trường mang tên Aogami với những chiếc đồng hồ quay ngược. Thế giới \hll\ ERROR đã chào đón họ bằng sương mù, tiếng cười vô hình và những bóng người tan biến giữa lớp học trống.

Nhưng câu chuyện không dừng lại ở nỗi sợ.

Xuyên suốt ba hồi chính, từ những bước chân dò dẫm đầy hoang mang trong hành lang vô tận của Épisode 0, đến việc lần lượt bóc tách từng lớp bí ẩn: vụ trường Aogami sụp đổ, những vụ mất tích bí ẩn quanh nhà ga, tiếng còi tàu ma quái trong đêm, và linh hồn báo thù của một cô gái luôn nở nụ cười hạnh phúc giữa khoảng không đổ nát - Kuon, Kaigou và Suzuki đã dần lấy lại danh tính thật sự của mình, dần hiểu ra rằng thế giới ERROR không đơn thuần là một chiều không gian bị lỗi, mà là một vũ trụ song song được dệt nên từ những mảnh vỡ linh hồn và tiềm thức mạnh mẽ nhất.

Ở trung tâm của tất cả là Misora Shino - cô gái nhút nhát, tự ti, người vừa là nạn nhân, vừa là nguồn gốc của mọi vòng lặp. Và bên cạnh cô là Shinomiya Sakura, người canh giữ đền Aogami, với những bí mật chôn vùi dưới tán cây ngàn năm. Khi sự thật được phơi bày, khi vòng lặp bị phá vỡ, thị trấn Aogami đã được tái sinh. Nhưng điều kỳ diệu nhất không nằm ở việc thế giới được ``sửa lại'', mà nằm ở việc một linh hồn bị lãng quên từ thuở xa xưa - Naomi - đã được trao cho một cuộc sống mới, một cơ hội để được yêu thương thật sự.

Không còn những trận chiến sống còn, không còn sương mù và bóng tối. Thay vào đó là ánh nắng ban mai chiếu qua ô cửa sổ lớp học, là tiếng cười đùa của hai mươi ba học sinh lớp 1-C cùng những vị khách đến từ chiều không gian khác. Là những buổi tập thể dục buổi sáng, những giờ ra chơi ồn ào, những bữa cơm trưa ấm cúng, những cuộc phiêu lưu nhỏ bé quanh thị trấn, và cả những đêm ngắm sao trên sân thượng. Đó là tuổi thanh xuân mà Kuon, Kaigou, Suzuki và Naomi chưa từng có được - một mùa hè trọn vẹn giữa vòng tay của những người bạn chân thành nhất.

Buổi tiệc ngủ cuối cùng với bảy hoạt động điên rồ, cuộc chiến gối lông chim tung tóe, bản nhạc ``Mùa hè cuối cùng ở Aogami'' của Yuki và Azuki, bản ``báo cáo thống kê cảm xúc'' đầy hài hước mà cũng đầy xúc động của Hệ thống Game, bức tranh của Mari, cuốn album ảnh của Uzuki, phong thư đóng sáp đỏ của cả lớp, bộ đồng phục và những bộ yukata được trao lại như kỷ vật - tất cả đều là những mảnh ghép nhỏ bé nhưng vô giá, góp phần tạo nên một bức tranh tuổi trẻ rực rỡ nhất mà tôi từng được viết ra.

BookER không phải là câu chuyện về những anh hùng cứu thế giới. Đó là câu chuyện về những con người bình thường, bị ném vào một hoàn cảnh phi thường, và lựa chọn sống trọn vẹn từng khoảnh khắc thay vì chìm đắm trong tuyệt vọng. Là câu chuyện về việc ``tìm lại mình'' giữa những mảnh vỡ của ký ức và danh tính. Là câu chuyện về tình bạn vượt qua ranh giới của không gian và thời gian.

Và trên hết, đó là câu chuyện về một câu hỏi giản đơn nhưng chứa đựng cả một bầu trời cảm xúc:

\vspace{1em}
\begin{center}
    \textit{``Các cậu sẽ mãi là bạn của mình chứ?''}
\end{center}
\vspace{1em}

Câu trả lời, tôi tin, nằm ở chính trái tim của mỗi người đã đồng hành cùng nhà Hikawa từ trang đầu tiên cho đến trang cuối cùng.

Và bây giờ, chúng ta hãy cùng quay trở lại cuộc sống thường nhật ở văn phòng \hll\ cùng với Hikawa Kuon. Những ngày mùa thu sắp tới cùng với những câu chuyện điên rồ hơn, những nhân vật mới đang chờ đón cậu ở phía trước.

Nhưng lần này, cậu sẽ không còn cô độc nữa.

\end{document}